\section{存在问题与未来趋势}
尽管深度学习在智慧能源中展现出明显潜力，但其工程化应用仍面临若干关键问题。首先是数据问题。智慧能源系统中的数据往往来自不同设备、不同采样频率和不同协议，存在缺失值、异常值、噪声和时间对齐误差。如果数据预处理不到位，再复杂的模型也难以稳定发挥作用。尤其在光伏预测中，云量突变、局地遮挡和传感器故障都会显著影响模型泛化。

其次是跨场景泛化问题。很多论文在单一站点上可以取得很好的结果，但换到另一地区、另一设备或另一季节后，误差会明显增大。这说明模型可能学到的是“站点特征”，而不是更一般的能源规律。元学习、迁移学习和跨站点集成是解决这一问题的重要方向，但目前仍处于快速发展阶段。

再次是可解释性问题。能源系统属于典型的高可靠性场景，调度人员通常不会仅凭“模型误差较小”就完全信任一个黑箱系统。未来如果要将深度学习更大规模地用于调度与控制，就必须给出更明确的解释机制，例如特征贡献分析、置信区间预测、物理边界校验和异常检测联动等。

最后是预测与决策之间的脱节问题。许多研究把预测误差作为唯一目标，但真实系统更关心的是调度成本、备用容量、安全约束和用户侧体验。未来智慧能源中的高价值方向，很可能不是孤立地追求更小的 RMSE，而是把预测、优化和控制纳入统一框架中，实现“预测-决策-反馈”闭环。
